This study develops deep learning surrogate models for the direct computation of Credit Valuation Adjustment (CVA) and Debit Valuation Adjustment (DVA) for an interest rate swap (IRS) under the Hull-White one-factor model. Standard xVA computation requires a nested Monte Carlo simulation in which the outer loop evolves the short rate along simulated paths and the inner loop re-prices the derivative at each monitoring date and scenario. This nested structure grows rapidly in computational cost with the number of scenarios and monitoring dates, and hence creates a computational bottleneck. We address this by training two differential deep learning surrogates to directly approximate the Expected Positive Exposure (EPE) and Expected Negative Exposure (ENE) profiles as explicit functions of the initial yield curve, Hull-White parameters, and monitoring time. Since EPE and ENE are already expectations over the short rate distribution, the short rate itself does not appear as a surrogate input, and CVA and DVA are recovered at inference time through a single forward pass per monitoring date with no simulation. Each surrogate is trained jointly on exposure values and their analytic curve sensitivities, which regularizes the learned function and improves accuracy under yield curve shapes not well represented in the training distribution. The trained surrogates reproduce EPE, ENE, CVA, and DVA estimates consistent with a full Monte Carlo benchmark across a range of yield-curve environments, whilst eliminating simulation cost during inference. These results establish that the deep learning surrogates offer a practically viable approach to xVA computation, and the framework is sufficiently general to extend to instruments with more complex payoff structures.
